% !TEX root = ../../cookbook.tex
\newpage
\recipe{Madeleines}{Pour 24 madeleines.}{}{}

\subsection*{Ingrédients}
\begin{ingredients}
	\item 175 g de farine
	\item 175 g de beurre
	\item 3 \oe ufs
	\item 150 g de sucre en poudre
	\item \sfrac{1}{2} \textit{càc} de levure chimique
	\item 2 pincées de sel
	\item Zeste d'une orange
	\item Zeste d'un citron
	\item 10 g de miel
	\item 1 \textit{càc} d'extrait de vanille
\end{ingredients}

\medskip

% --- Preparation ---
Faire fondre le \textbf{beurre} au micro-ondes et laisser tiédir.
Dans un bol, bien mélanger \textbf{tous les poivres} et les \textbf{zestes}.
Ajouter le \textbf{beurre} fondu, mélanger, puis ajouter les \textbf{\oe ufs} et la \textbf{vanille}.
Bien mélanger\footnote{Mais ne pas battre, sinon les madeleines seront trop gonflées et légères.} jusqu'à obtenir une crème lisse et homogène.

Laisser reposer au réfrigérateur pendant au moins 1 h.

Préchauffer le four à 180°C, beurrer un ou plusieurs moules à madeleines, et y répartir la pâte dans 24 empreintes.

Cuire pendant 13 minutes.
Les madeleines sont prêtes lorsqu'elles sont dorées et qu'elles ont formé leur bosse caractéristique.


\vspace*{0.75 cm}

% --- Description ---
\begin{recipeblock}
	\noindent 
\end{recipeblock}